\chapter{Geltungsbereich und Strukturbegriffe}
\label{b45guide:ch:foundations}

Dieser Leitfaden ergänzt Band~45 um eine Auswahl wichtiger Definitionen,
Charakterisierungen und Folgerungen für \emph{Halbverbände und Verbände}.
Er enthält ausschließlich Aussagen und Begriffsbestimmungen, keine
Beweise oder Beweisskizzen. Die Nummern mit Präfix $45H$ gehören
zu dieser Ergänzung. Der Leitfaden beansprucht keine vollständige
Aufzählung aller bekannten oder ableitbaren Resultate.

\section{Voraussetzungen und Bandgrenzen}

Die Verknüpfungen sind überall definierte binäre Operationen auf dem
angegebenen Träger. Wie im Hauptband sind leere Träger zugelassen,
solange keine Konstanten wie $0$ oder $1$ verlangt werden.
Endlichkeit und das Vorhandensein solcher Elemente werden jeweils
ausdrücklich vorausgesetzt. Als mengentheoretischer Rahmen dient ZFC.

Die Auswahl betrifft die Halbverbands- und Verbandsoperationen:
ihre Axiomatik, Erzeugung, Darstellungen und operationserhaltende
Abbildungen. Ein Satz wird nicht allein deshalb aufgenommen,
weil ein allgemeines Ergebnis über Funktionen oder Ordnungen
auch auf diese Strukturen anwendbar ist.

\begin{center}
\small
\begin{tabularx}{\linewidth}{@{}p{.72\linewidth}X@{}}
\toprule
Eigenständiges Thema außerhalb dieses Leitfadens & Band\\
\midrule
Allgemeine Funktionen; Injektivität, Surjektivität, Bijektivität & 5--8\\
Allgemeine Äquivalenzrelationen & 11\\
Allgemeine Halbordnungen & 12\\
Allgemeine Schranken-, Infimums- und Supremumssätze & 13\\
Rechenregeln für vorhandene Paarinfima und Paarsuprema & 14\\
Totale Ordnungen und Wohlordnungen & 15--16\\
Allgemeine Theorie endlicher Mengen & 20\\
Andere algebraische Strukturtheorien & 28--44\\
Anwendungen auf Frankls Vermutung & 46\\
\bottomrule
\end{tabularx}
\end{center}

Distributive, modulare, komplementierte und vollständige Verbände
werden nicht als zusätzliche axiomatische Klassen entwickelt.
Freie Objekte, Unterstrukturen, Ideale und Kongruenzen erscheinen
nur als Konstruktionen innerhalb der Halbverbands- und Verbandstheorie.
Bei Kongruenzen ist ausschließlich die Verträglichkeit mit den
betrachteten Operationen Gegenstand; allgemeine Aussagen über
Äquivalenzrelationen werden vorausgesetzt.

\section{Algebraische Definitionen}

\begin{definition}[Halbverband]
\label{b45guide:def:semilattice}
Ein \emph{Halbverband} oder \emph{Halbgitter} ist ein Paar $(S,\star)$
aus einer Menge $S$ und einer Operation $\star:S\times S\to S$ mit
\[
\begin{aligned}
 (x\star y)\star z&=x\star(y\star z),&
 x\star y&=y\star x,& x\star x&=x
 \qquad(x,y,z\in S).
\end{aligned}
\]
Dies sind Assoziativität, Kommutativität und Idempotenz.
Im Sinne des Hauptbands schreiben wir $\Semi{S,\star}$.
\end{definition}

\begin{definition}[Verband]
\label{b45guide:def:lattice}
Ein \emph{Verband} ist eine Struktur $(L,\wedge,\vee)$,
in der $(L,\wedge)$ und $(L,\vee)$ Halbverbände sind und
für alle $x,y\in L$ die beiden Absorptionsgesetze gelten:
\[
 x\wedge(x\vee y)=x,\qquad x\vee(x\wedge y)=x.
\]
Die Bezeichnung im Hauptband lautet $\LatticeAlg{L,\wedge,\vee}$.
Es werden keine weiteren Identitäten vorausgesetzt.
\end{definition}

\begin{definition}[Induzierte Vergleiche und duale Lesart]
\label{b45guide:def:order}
Eine Halbverbandsoperation kann supremal oder infimal gelesen werden:
\[
 x\leq_{\vee}y\iff x\vee y=y,
 \qquad
 x\leq_{\wedge}y\iff x\wedge y=x.
\]
Die Zeichen $\vee$ und $\wedge$ kennzeichnen hier die jeweilige
Lesart. Für dieselbe Operation $\star$ sind die beiden Relationen
zueinander umgekehrt. In einem Verband verwenden wir $\leq$
für den gemeinsamen Vergleich; $x<y$ bedeutet $x\leq y$ und $x\ne y$.
\end{definition}

\begin{theorem}[Algebraische Charakterisierung der Verbandsstruktur]
\label{b45guide:thm:compatible}
Seien $\wedge,\vee$ zwei Halbverbandsoperationen auf derselben Menge $L$.
Dann sind äquivalent:
\begin{enumerate}
\item $(L,\wedge,\vee)$ ist ein Verband.
\item Für alle $x,y\in L$ gilt
\[
 x\wedge y=x\quad\Longleftrightarrow\quad x\vee y=y.
\]
\end{enumerate}
Die gemeinsame Relation ist eine Halbordnung und bestimmt beide
Operationen eindeutig. Insbesondere gibt es zu einer vorgegebenen
Halbverbandsoperation höchstens eine zweite Operation, die mit ihr
in dieser Rollenverteilung einen Verband bildet. \cite{project,nation}
\end{theorem}

Die Charakterisierung benennt die Verbindung zur Ordnung.
Die allgemeine Existenz- und Rechentheorie der zugehörigen
Paarinfima und Paarsuprema bleibt in den Bänden~13 und~14.

\begin{theorem}[Dualitätsprinzip]
\label{b45guide:thm:duality}
Ist $(L,\wedge,\vee)$ ein Verband, so auch $(L,\vee,\wedge)$.
Seine induzierte Ordnung ist die umgekehrte Ordnung von $L$.
Jede für alle Verbände gültige Aussage bleibt gültig, wenn
$\wedge$ und $\vee$ sowie $\leq$ und $\geq$ überall vertauscht werden.
Bei mitgeführten Schranken werden zusätzlich $0$ und $1$ vertauscht.
\end{theorem}

\begin{definition}[Null, Eins und beschränkte Strukturen]
\label{b45guide:def:bounds}
In einem supremal gelesenen Halbverband heißt $0\in S$ eine
\emph{Null}, wenn $0\vee x=x$ für alle $x\in S$ gilt.
Eine \emph{Eins} ist ein Element $1\in S$ mit $x\vee1=1$.
Ein \emph{beschränkter Verband} $(L,\wedge,\vee,0,1)$ ist ein
Verband mit
\[
 0\vee x=x,\quad 0\wedge x=0,\quad
 1\vee x=1,\quad 1\wedge x=x .
\]
Die Schranken sind, falls vorhanden, eindeutig. Eine benannte Null
oder Eins wird nur dort als Teil der Struktur mitgeführt, wo dies
ausdrücklich angegeben ist.
\end{definition}
